
\section{Threats to Validity and Limitations}
\label{sec:limitations}

\paragraph{Synthetic formal benchmark.}
The 180 scenarios are constructed from formal world specifications. This enables exact
paired interventions, dual validation, and contract-level supportedness claims, but not
human-preference or ecological-validity claims.

\paragraph{Deterministic policy.}
The confirmatory study uses one frozen deterministic policy. This isolates the
evidence-surface mechanism but does not establish prevalence across LLM agents or other
decision policies.

\paragraph{Family and policy construction.}
The direct C1 effect is concentrated entirely in F6. The frozen policy was intentionally
designed with explicit visible-gap fallbacks for several other obligation types, so this
heterogeneity is partly a property of the policy--benchmark pair. The result identifies a
controlled mechanism; it does not estimate how often each failure type occurs in deployed
agents.

\paragraph{Formal obligation oracle.}
C5 receives an accurate completeness signal derived from the benchmark's formal
obligations. Producing such a signal robustly in open-world systems remains unresolved.

\paragraph{Retriever realism.}
C2 is a lexical baseline; C3 is a deterministic hash-vector baseline rather than a
pretrained semantic encoder; C4 is their fixed reciprocal-rank fusion. These are controlled
retrieval mechanisms, not claims about state-of-the-art retrieval. All use frozen
top-$k=2$; different budgets may alter coverage and failure patterns.

\paragraph{Inferential scope.}
McNemar and bootstrap results summarize the paired 180-scenario benchmark under its frozen
construction. The scenarios are not a probability sample from a defined real-world
population, so the p-value and interval must not be interpreted as population-prevalence
estimates.

\paragraph{Coverage-signal errors are not studied.}
C5 assumes that the signal itself is correct. False-complete coverage is likely the most
safety-relevant next failure mode.

\paragraph{No deployment claim.}
The results do not establish safety in real decision-making systems or high-stakes domains.
